\documentclass[11pt,a4paper]{article}

% --- Mathematics (must precede unicode-math) ---
\usepackage{amsmath}
\usepackage{mathtools}

% --- Fonts (lualatex) ---
\usepackage{fontspec}
\usepackage{unicode-math}
\setmathfont{KpMath-Regular.otf}
\setmathfont{KpMath-Bold.otf}[version=bold]

% --- Microtypography (after all font setup) ---
\usepackage{microtype}

% --- Colors & page layout ---
\usepackage[dvipsnames]{xcolor}
\usepackage[margin=25mm, top=30mm, bottom=30mm]{geometry}

% --- Headers & footers ---
\usepackage{fancyhdr}
\pagestyle{fancy}
\fancyhf{}
\fancyhead[L]{\small\itshape Mode-Converted Reflections and Evanescent Incidence}
\fancyhead[R]{\small\itshape Nestor (2026)}
\fancyfoot[C]{\thepage}
\renewcommand{\headrulewidth}{0.4pt}
\renewcommand{\footrulewidth}{0pt}

% --- Section numbering ---
\setcounter{secnumdepth}{3}

% --- Tables ---
\usepackage{booktabs}
\usepackage{array}

% --- Captions ---
\usepackage[font=small, labelfont=bf, labelsep=period]{caption}

% --- Spacing ---
\usepackage{setspace}
\setstretch{1.15}
\setlength{\parskip}{0.4em}

% --- Hyperlinks (last before macros) ---
\usepackage[colorlinks=true, linkcolor=blue!60!black,
            citecolor=blue!60!black, urlcolor=blue!60!black]{hyperref}

% --- Macros (house set) ---
\newcommand{\bG}{\mathbf{G}}
\newcommand{\bI}{\mathbf{I}}
\newcommand{\bx}{\mathbf{x}}
\newcommand{\dd}{\mathrm{d}}
\newcommand{\order}[1]{\mathcal{O}\!\left(#1\right)}
\newcommand{\bT}{\mathbf{T}}
\newcommand{\Dlam}{\Delta\lambda}
\newcommand{\Dmu}{\Delta\mu}
\newcommand{\Drho}{\Delta\rho}
\newcommand{\re}{\operatorname{Re}}
\newcommand{\im}{\operatorname{Im}}

% --- New macros for this document ---
\newcommand{\bu}{\mathbf{u}}
\newcommand{\bfv}{\mathbf{f}}
\newcommand{\bsv}{\mathbf{s}}
\newcommand{\bd}{\hat{\mathbf{d}}}
\newcommand{\bsig}{\boldsymbol{\sigma}}
\newcommand{\rhat}{\hat{\mathbf{r}}}
\newcommand{\svhat}{\hat{\mathbf{e}}_{\mathrm{SV}}}
\newcommand{\shhat}{\hat{\mathbf{e}}_{\mathrm{SH}}}
\newcommand{\etaP}{\eta_{P}}
\newcommand{\etaS}{\eta_{S}}
\newcommand{\Rpsv}{\mathbf{R}}
\newcommand{\code}[1]{\texttt{#1}}

\begin{document}

% --- Title block ---
\thispagestyle{plain}
\begin{center}
\vspace*{1em}
{\LARGE\bfseries Mode-Converted Specular Reflections\\[0.3em]
and Evanescent Incidence in the Periodic Slab Solver}\\[1.8em]
{\large T.\,M.\ Nestor}\\[0.5em]
{\itshape Research School of Earth Sciences, Australian National University}\\[2em]
\end{center}

% =============================================================================
\section{Introduction and scope}
\label{sec:intro}
% =============================================================================

This document records three connected developments, each validated
against an independent exact reference.

\begin{enumerate}
\item \textbf{Five-channel Mie verification (sphere).}  The elastic Mie
  solution and the Foldy--Lax voxelised sphere were compared
  phase-sensitively (real and imaginary parts separately) across all
  five non-zero channels of the $3\times 3$ scattering matrix
  ($P{\to}P$, $P{\to}SV$, $SV{\to}P$, $SV{\to}SV$, $SH{\to}SH$) at
  $ka = 0.1$, $0.5$, and $1.5$.
\item \textbf{Full specular reflection matrix of the periodic
  heterogeneous slab.}  The Weyl lattice-sum extraction, previously
  limited to $R_{PP}$, was generalised to the $2\times 2$ P--SV matrix
  plus the scalar SH channel at arbitrary horizontal slowness~$p$, and
  validated channel-by-channel against Kennett reflectivity.  The
  ocean-bottom study now propagates the full $2\times 2$ sub-ocean
  reflectivity, so that internal $P{\to}S{\to}P$ conversion contributes
  to the observable water-column $R_{PP}$.
\item \textbf{Physical evanescent incidence.}  Incident fields are
  constructed from complex slowness vectors, so post-critical incidence
  ($p > 1/\alpha$) is represented by the true exponentially decaying
  inhomogeneous wave rather than a homogeneous-wave surrogate.
\end{enumerate}

Two pre-existing defects were found and repaired along the way
(Section~\ref{sec:bugs}), and two normalisation conventions were pinned
by measurement rather than assumption
(Sections~\ref{sec:sphere-renorm} and~\ref{sec:convention}).  All
numerical values quoted below are reproduced by the committed test
suite (641 tests at the time of writing).

Throughout, coordinates follow the package convention: $z$ is axis~0
(positive down), $x$ is axis~1, $y$ is axis~2; the sagittal plane is
$z$--$x$.  Vertical slownesses are
$\eta_m = \sqrt{1/c_m^2 - p^2}$ with the branch $\im\eta_m \ge 0$, and
all complex polarisation algebra uses analytic continuation (no
conjugation).

% =============================================================================
\section{Five-channel Mie verification for the sphere}
\label{sec:sphere}
% =============================================================================

For an isotropic sphere the P--SV and SH problems decouple at each
partial-wave order, leaving five non-zero channels of the scattering
matrix.  $P$-wave incidence is axially symmetric ($m=0$); $SV$ and $SH$
plane-wave incidence carry azimuthal order $m=1$, with far-field angular
functions $P_n^1(\cos\theta) = \dd P_n/\dd\theta$ (Condon--Shortley),
$\dd P_n^1/\dd\theta$, and $P_n^1/\sin\theta$ respectively.

\subsection{The $m{=}1$ renormalisation}
\label{sec:sphere-renorm}

The scattering coefficients are stored in the $m{=}0$ plane-wave
expansion normalisation.  Evaluating $SV$/$SH$ far fields requires the
$m{=}1$ expansion coefficient of $\hat{\mathbf{x}}\,e^{ik_S z}$, which
differs by the measured factor
\begin{equation}
  \mathrm{renorm}(n) \;=\; -\,\frac{1}{n(n+1)},
  \label{eq:renorm}
\end{equation}
the minus sign arising from the
$\nabla\times\nabla\times(\mathbf{r}\,\psi)$ convention.  (The planning
estimate $+i/n(n+1)$ was wrong; the correct factor was pinned by the
Rayleigh-limit comparison and confirmed by reciprocity.)  A further
empirical subtlety: the $SH{\to}SH$ far field at $\varphi = 0$ is
carried by the N-type (poloidal, $b_n^{SV}$) coefficients; the true
M-type ($c_n$) contribution is $\order{(ka)^4}$ and negligible.

\subsection{Reciprocity and results}

The off-diagonal channels obey
\begin{equation}
  k_P^2\, f_{PS}(\theta) \;=\; -\,k_S^2\, f_{SP}(\theta),
  \label{eq:sphere-recip}
\end{equation}
with the conversion constant exactly $-1$ (measured ratio
$-1.0002 + 0.0002i$ at $ka = 0.05$, $-1.0000$ at $ka = 0.5$; asserted in
\code{TestReciprocity}).  Phase-sensitive errors between Mie and the
voxelised Foldy--Lax sphere (relative to the pattern peak; $\im$ errors
were ${\approx}\,0$ in every case):

\begin{center}
\begin{tabular}{lccc}
\toprule
Channel & $ka=0.1$ ($n_{\rm sub}{=}4$) & $ka=0.5$ ($n_{\rm sub}{=}6$)
        & $ka=1.5$ ($n_{\rm sub}{=}8$) \\
\midrule
$P{\to}P$   & 0.041 & 0.200 & 0.316 \\
$P{\to}SV$  & 0.038 & 0.202 & 0.393 \\
$SV{\to}P$  & 0.048 & 0.186 & 0.368 \\
$SV{\to}SV$ & 0.047 & 0.228 & 0.565 \\
$SH{\to}SH$ & 0.047 & 0.229 & 0.573 \\
\bottomrule
\end{tabular}
\end{center}

The errors scale as $\order{ka}$ (voxelisation and T-matrix truncation
at finite contrast), falling consistently under mesh refinement.  The
identical Mathematica pipeline (\code{FiveChannelExtension.wl}, run
headlessly by \code{FiveChannelDriver.wl}) reproduces all five channels
to better than 10\% at $k_S a = 0.52$ with $9$ voxels per diameter, with
the $SV$/$SH$ channel errors agreeing to ten digits --- a free
symmetry check of the two independent solves.

% =============================================================================
\section{Specular reflection matrix of the periodic slab}
\label{sec:slab}
% =============================================================================

\subsection{Generalised Weyl extraction}

For the infinite periodic slab the 2D lattice sum replaces the
spherical-wave kernel by its specular Weyl term,
$e^{ikr}/4\pi r \to i\,e^{i\omega\eta_m |z|}/(2\,\omega\eta_m d^2)$ per
mode.  With horizontally layer-averaged Foldy--Lax sources
$(\bfv_\ell, \bsig_\ell)$ and the upgoing slowness vector
$\bsv_m = (-\eta_m,\, p,\, 0)$, complex-unit direction
$\bd_m = c_m \bsv_m$, the outgoing amplitudes are
\begin{align}
  Q_{P,\ell} &= -\,\bd_P\!\cdot\!\bfv_\ell
               \;-\; i\omega\,\bsv_P\!\cdot\!\bsig_\ell\!\cdot\!\bd_P,
  &
  R_P &= \frac{-i}{2\,\omega\etaP\, d^2 \rho\,\alpha^2}
         \sum_\ell Q_{P,\ell}\; e^{\,i\omega\etaP z_\ell},
  \label{eq:weylP}\\
  \mathbf{Q}_{S,\ell} &= -\,\bfv_\ell
               \;-\; i\omega\,\bsig_\ell\!\cdot\!\bsv_S,
  &
  R_{SV,SH} &= \frac{-i}{2\,\omega\etaS\, d^2 \rho\,\beta^2}
         \sum_\ell
         \bigl(\svhat,\shhat\bigr)\!\cdot\!\mathbf{Q}_{S,\ell}\;
         e^{\,i\omega\etaS z_\ell},
  \label{eq:weylS}
\end{align}
with $\svhat = \beta\,(p,\,\etaS,\,0)$ and
$\shhat = (0,\,0,\,1)$.  Only the force term is negated (the T-matrix
convention $+V\omega^2\Drho^* \bu$ is opposite to the
Lippmann--Schwinger body force).  The full reflection matrix is
assembled from three periodic solves (incident $P$, $SV$, $SH$ at the
same slowness), sharing the local T-matrices.

\subsection{Two pre-existing defects}
\label{sec:bugs}

\textbf{Oblique stress coupling.}  The original $R_{PP}$ extraction used
the vertical wavenumber in the stress term,
$-i k_z\,(\bd\!\cdot\!\bsig\!\cdot\!\bd)$ with $k_z = \omega\etaP$.  The
correct coupling is the full wave vector,
$-i\omega\,(\bsv\!\cdot\!\bsig\!\cdot\!\bd) =
 -i k_P\,(\bd\!\cdot\!\bsig\!\cdot\!\bd)$.  The two coincide only at
$p=0$ --- and no oblique slab-versus-Kennett test existed, so the defect
had never been observed.  Equation~\eqref{eq:weylP} is the corrected
form; the oblique validation of Table~\ref{tab:channels} is its proof.

\textbf{$p\to 0$ polarisation discontinuity.}  The incident-field
builder's vertical-incidence special case selected the $SV$ polarisation
$+\hat{\mathbf{x}}$, whereas the $p\to 0^{+}$ limit of its own generic
branch is $-\hat{\mathbf{x}}$.  The discontinuity inverted the sign of
$R_{SS}$ at exactly $p=0$, made observable by Kennett's convention
$K_{SS}(0) = -K_{SH}(0)$ (the $SV$ sign flips under reflection; same
physics, opposite bookkeeping to $SH$).  The special case now takes the
continuous limit.

\subsection{The convention bridge to Kennett}
\label{sec:convention}

The slab extraction produces displacement-amplitude ratios; the package's
Kennett implementation stores reflection matrices in its modified
(eigenvector-normalised) convention, in which they are symmetric.  The
planned conversion $\tilde{R}_{ij} = \sqrt{\eta_i/\eta_j}\,R_{ij}$
proved wrong by measurement.  The correct bridge is the diagonal
similarity
\begin{equation}
  \tilde{\Rpsv} \;=\; \mathbf{D}\,\Rpsv\,\mathbf{D}^{-1},
  \qquad
  \mathbf{D} = \operatorname{diag}\!\bigl(
     \alpha\sqrt{\etaP},\;\; i\,\beta\sqrt{\etaS}\bigr),
  \label{eq:similarity}
\end{equation}
the velocity factors and the factor $i$ on $SV$ reflecting Kennett's
eigenvector normalisation (the $-2i$ factor in the off-diagonal
construction of \code{psv\_solid\_solid}).  The decisive,
convention-independent check is the reciprocity invariant: the product
of the two off-diagonal slab/Kennett ratios is invariant under
\emph{any} diagonal similarity and measured $0.987 \approx 1$, proving
the physics correct before the bookkeeping was; the per-channel ratios
under~\eqref{eq:similarity} are then real and within 1\% of unity.
Diagonal entries are unchanged by any diagonal similarity, so the
historical $R_{PP}$ results are unaffected.

\subsection{Channel validation}

\begin{table}[ht]
\centering
\caption{Uniform slab versus Kennett (modified convention).  Background
$\alpha = 5000$, $\beta = 3000$\,m/s, $\rho = 2500$\,kg/m$^3$; contrast
$\Dlam = 2$, $\Dmu = 1$\,GPa, $\Drho = 100$\,kg/m$^3$; $\omega = 150$
rad/s; $H = 4$\,m, mesh $a = 2$\,m, $N_z = 1$ unless stated.  Relative
errors.}
\label{tab:channels}
\begin{tabular}{lccc}
\toprule
Channel & $p = 0$ & $p = 1{\times}10^{-4}$\,s/m
        & $p = 2.5{\times}10^{-4}$\,s/m (post-critical) \\
\midrule
$R_{PP}$ & 0.0045 & 0.0048 & 0.054 \quad(0.0009 at $a{=}0.5$, $N_z{=}4$) \\
$R_{PS}$ & (vanishes; $<10^{-2}$ rel.) & 0.0066 & 0.033 \\
$R_{SP}$ & (vanishes) & 0.0066 & 0.033 \\
$R_{SS}$ & 0.0032 & 0.0110 & 0.043 \\
$R_{SH}$ & 0.0032 & 0.0057 & 0.044 \\
\bottomrule
\end{tabular}
\end{table}

The modified matrix is symmetric to $1.3\times 10^{-5}$ at oblique
incidence (independent solves: $R_{SV}$ from $P$ incidence versus $R_P$
from $SV$ incidence), the conversion channel is Born-linear at weak
contrast (doubling ratio $2.0$, rtol $0.05$), and at $p = 0$ the
degeneracy $R_{SH} = -\tilde{R}_{SS}$ holds at machine level (the $SV$
and $SH$ systems on the square lattice are exact $90^\circ$ rotations of
one another).

% =============================================================================
\section{Ocean-bottom embedding: the
  \texorpdfstring{$2\times 2$}{2x2} sub-ocean recursion}
\label{sec:ocean}
% =============================================================================

The ocean-bottom study (water $|$ heterogeneous sediment slab $|$
elastic halfspace) previously injected only the scalar $R_{PP}$ of the
slab into a scalar water-step recursion.  Internal mode conversion ---
$P$ converting to $SV$ inside the sediment package, reverberating, and
converting back --- was absent from the observable.  The sub-ocean
reflectivity is now the full matrix
\begin{equation}
  \mathbf{M}(\omega)
  \;=\;
  \mathbf{E}\,\mathbf{R}^{\rm sed|hs}\,\mathbf{E}
  \;+\;
  \tilde{\Rpsv}^{\rm slab}(\omega),
  \qquad
  \mathbf{E} = \operatorname{diag}\!\bigl(e^{i\omega\etaP^{\rm sed}H},\,
                                          e^{i\omega\etaS^{\rm sed}H}\bigr),
\end{equation}
so a conversion path acquires the mode-correct two-way phase
$E_i E_j$, and the water-side observable is
\begin{equation}
  R_{PP}^{\rm obs}
  \;=\;
  \Bigl[\mathbf{R}_d + \mathbf{T}_u\,\mathbf{M}\,
        \bigl(\bI - \mathbf{R}_u\mathbf{M}\bigr)^{-1}\mathbf{T}_d
  \Bigr]_{00},
\end{equation}
with the fluid--solid coefficients of \code{psv\_fluid\_solid}.  $SH$
cannot couple through the fluid and is omitted (the slab-level $R_{SH}$
remains available directly).  Per-frequency cost is two Foldy--Lax
solves ($P$- and $SV$-incident).

\textbf{Anchors.}  At zero contrast the decomposed path reproduces the
full-stack Kennett $R_{PP}$ to relative error $2\times 10^{-16}$ at both
$p = 0$ and $p = 2\times 10^{-4}$\,s/m (asserted at rtol $10^{-10}$).
The oblique case is the load-bearing anchor: at $p=0$ every quantity
introduced by the matrix upgrade is invisible, and a deliberately
injected phase error of the $E_iE_j$ type produces a 133\% error that
only the oblique anchor detects.

\textbf{Magnitude of the restored physics.}  At moderate contrast and
oblique incidence, suppressing the off-diagonal entries of $\mathbf{M}$
shifts the observable $R_{PP}$ at the $\sim$10\% level --- dominated by
the background sediment/halfspace interface conversion, which the old
scalar recursion also lacked.  The slab's own conversion contribution is
second order in the contrast, isolated by subtracting only
$\tilde{\Rpsv}^{\rm slab}$ off-diagonals: $1.2\times 10^{-5}$ relative
at weak contrast.  (An instructive trap: zeroing the off-diagonals of
the \emph{total} $\mathbf{M}$ measures the contrast-independent
background conversion, not the slab's --- the committed tests separate
the two.)

% =============================================================================
\section{Evanescent incidence}
\label{sec:evanescent}
% =============================================================================

Past the critical slowness ($p > 1/\alpha$ for $P$) the incident field
is an inhomogeneous wave: with $\etaP = i|\etaP|$,
\begin{equation}
  \bu^{\rm inc}(\bx)
  \;=\;
  \mathbf{p}\; e^{\,i\omega\,\bsv\cdot\bx}
  \;=\;
  \mathbf{p}\; e^{\,i\omega p x}\, e^{-\omega|\etaP| z},
  \qquad
  \bsv = (\etaP,\, p,\, 0),\quad \mathbf{p} = \alpha\,\bsv ,
\end{equation}
decaying with depth.  The incident-field builder originally degenerated
to a horizontally propagating homogeneous wave there (constant amplitude
at all depths), and the $P$-incident column of the reflection matrix
carried a documented warning.  The builder now constructs fields
directly from complex slowness vectors and polarisations
($\mathbf{p}\cdot\mathbf{p} = 1$ by analytic continuation --- the same
convention as Kennett, though the Euclidean amplitude exceeds unity
post-critically); strains follow as
$\varepsilon_{ij} = \tfrac{i\omega}{2}(s_i p_j + s_j p_i)$.  The
per-layer decay ratio matches $e^{-\omega|\etaP|\,\Delta z}$ to
$10^{-10}$, and the sub-critical path reproduces the previous builder to
$2.6\times 10^{-16}$.

\begin{center}
\begin{tabular}{lcc}
\toprule
Channel ($p = 2.5\times 10^{-4}$\,s/m) & homogeneous surrogate & evanescent field \\
\midrule
$R_{PP}$ & 0.637 & 0.054 \\
$R_{PS}$ & 0.965 & 0.033 \\
$R_{SP}$ & 0.033 & 0.033 \\
\bottomrule
\end{tabular}
\end{center}

Accuracy is now symmetric across the critical slowness (errors 0.025 at
$p = 1.9\times 10^{-4}$ versus 0.031 at $2.1\times 10^{-4}$), and the
residual post-critical error is pure vertical discretisation: a
monotonic exponential sampled by one-point quadrature biases the layer
average (unlike oscillatory sub-critical phase, which self-cancels), and
the error collapses $0.054 \to 0.009 \to 0.0009$ under refinement
$(a, N_z) = (2,1) \to (1,2) \to (0.5,4)$ at fixed $H$ --- pinned by a
committed convergence test.  The conversion channels floor near 2\%,
the known Rayleigh T-matrix truncation scale at this contrast.  The only
remaining special points are the exact grazing slownesses
$p = 1/\alpha,\, 1/\beta$, where the Weyl prefactor $1/\eta$ is
genuinely singular.

% =============================================================================
\section{Code artifacts}
\label{sec:code}
% =============================================================================

\begin{center}
\small
\begin{tabular}{ll}
\toprule
Artifact & Role \\
\midrule
\code{sphere\_scattering.py} &
  five-channel Mie far fields, $SV$/$SH$ decomposition \\
\code{FiveChannelExtension.wl}, \code{FiveChannelDriver.wl} &
  Mathematica five-channel comparison (headless) \\
\code{slab\_weyl\_amplitudes} &
  shared $P$/$SV$/$SH$ Weyl extractor, eqs.~\eqref{eq:weylP}--\eqref{eq:weylS} \\
\code{SlabReflectionMatrix.to\_modified} &
  convention bridge, eq.~\eqref{eq:similarity} \\
\code{slab\_reflection\_matrix} &
  three-solve matrix assembly at slowness $p$ \\
\code{kennett\_reference\_matrix} &
  five-channel Kennett reference for a uniform layer \\
\code{\_build\_slab\_incident\_field\_slowness} &
  complex-slowness (evanescent) incident fields \\
\code{\_kennett\_water\_step} ($2\times 2$), \code{MT\_psv} &
  ocean-bottom sub-ocean matrix recursion \\
\bottomrule
\end{tabular}
\end{center}

All quoted measurements are asserted by the test suite
(\code{TestCompleteScatteringMatrix}, \code{TestSlabReflectionMatrix},
\code{TestEvanescentIncidence}, \code{TestModeConvertedOceanBottom},
and the oblique zero-contrast anchors), 641 tests passing at the time of
writing.

\end{document}
